\documentclass{article}

\usepackage[margin=1in]{geometry}
\usepackage{amsmath,amsthm,amssymb}
\usepackage[spanish]{babel}
\usepackage{enumitem}

\theoremstyle{plain}
\newtheorem{proposicion}{Proposición}[section]

\theoremstyle{definition}
\newtheorem{definition}{Definición}[section]

% Number sets
\newcommand{\R}{\mathbb{R}}
\newcommand{\Z}{\mathbb{Z}}
\newcommand{\N}{\mathbb{N}}
\newcommand{\Q}{\mathbb{Q}}
\newcommand{\C}{\mathbb{C}}

% Topology operators
\newcommand{\cl}[1]{\overline{#1}}              % Closure
\newcommand{\Int}[1]{{#1}^\circ}                % Interior
\newcommand{\bd}{\partial}                      % Boundary
\newcommand{\Ent}[1]{\mathcal{O}(#1)}           % Neighborhoods (Entornos)
\newcommand{\pf}{\mathcal{P}_F}                 % Finite parts
\newcommand{\partes}{\mathcal{P}}               % Power set

% Useful shortcuts
\newcommand{\sii}{\iff}                          % Si y sólo si
\newcommand{\imp}{\implies}                      % Implies
\newcommand{\setdiff}{\smallsetminus}            % Set difference

% Net convergence/accumulation notation
\newcommand{\evto}[1]{\stackrel{E}{\hookrightarrow} #1}   % Eventually in
\newcommand{\frto}[1]{\stackrel{F}{\hookrightarrow} #1}   % Frequently in
\newcommand{\nnevto}[1]{\not\!\stackrel{E}{\hookrightarrow} #1}
\newcommand{\nnfrto}[1]{\not\!\stackrel{F}{\hookrightarrow} #1}

% Exercise environment
\newenvironment{ejercicio}[1]{\begin{trivlist}
\item[\hskip \labelsep {\bfseries Ejercicio}\hskip \labelsep {\bfseries #1.}]}{\end{trivlist}}

\setlist[enumerate,1]{label=(\alph*), leftmargin=*, itemsep=2pt}

\begin{document}

\title{Topología: Práctica V}
\author{Bustos Jordi\\Funciones continuas, métricas y topologías}

\maketitle

\section*{Funciones Continuas}

\begin{ejercicio}{1}
    Sean \( (X, \tau) \), \( (Y, \sigma) \) dos espacios topológicos y sea \( f : X \to Y \) una función. Demostrar las siguientes proposiciones:
    \begin{enumerate}
        \item La función \( f \) es continua si y sólo si para toda red \( \mathbf{x} = {(x_i)}_{i \in I} \) en \( X \) que converge a algún \( x \in X \) se cumple que \( f(x_i) \xrightarrow{i \in I} f(x) \).
        \item La función \( f \) es continua de \( (X, \tau) \to (Y, \sigma) \) si y sólo si es continua de \( (X, \tau) \to (f(X), \sigma_{f(X)}) \).
        \item Si \( f \) es continua y suryectiva, y \( D \) es \( \tau \)-denso en \( X \) entonces \( f(D) \) es \( \sigma \)-denso en \( Y \).
        \item Si \( f \) es continua y \( D \) es \( \tau \)-denso en \( X \) entonces \( f(D) \) es \( \sigma_{f(X)} \)-denso en \( f(X) \).
        \item Si \( f \) es continua y \( A \subseteq X \) entonces \( f|_A : A \to Y \) es continua.
    \end{enumerate}
\end{ejercicio}
\begin{proof}
    Dados \( (X, \tau) \), \( (Y, \sigma) \) dos espacios topológicos y sea \( f : X \to Y \) una función.
    \begin{enumerate}
        \item \( (\Leftarrow) \) \textbf{Proposición 5.1.1} del apunte de Demetrio. \( (\Rightarrow) \) Dado \( x \in X \) y \( V \in \mathcal{O}^a(f(x)) \) vale que \( f^{-1}(V) \in \tau \) por la continuidad de \( f \),
              como \( f^{-1}(V) \in \mathcal{O}^a(x) \) vale que: \[
                  x \evto{f^{-1}(V)} \implies f(x) \evto{V} \implies f(x_i) \xrightarrow{i \in I} f(x)
              \]
        \item \( (\Rightarrow) \) como \( \sigma_{f(X)} \subseteq \sigma \) entonces la \( \sigma \)-convergencia de redes implica la \( \sigma_{f(X)} \)-convergencia de redes.
              De lo que se deduce que \( f \) es continua por el punto anterior. \( (\Leftarrow) \) Si \( f \) es continua en \( (f(X), \sigma_{f(X)}) \) entonces \[
                  \forall U \in \sigma_{f(X)} \quad U = V \cap f(X) \text{ con } V \in \sigma \implies f^{-1}(U) = f^{-1}(V) \in \tau
              \]
              pero entonces \( f \) es continua de \( (X, \tau) \to (Y, \sigma) \).
        \item Dado \( D \) \( \tau \)-denso en \( X \) y \( f \) continua y suryectiva, dado \( y \in Y \) existe \( x \in X \) tal que \( f(x) = y \).
              Como \( D \) es denso en \( X \) entonces existe una red \( \mathbf{d} = {(d_i)}_{i \in I} \) en \( D \) que converge a \( x \) y por la continuidad de \( f \) se deduce que \( f(d_i) \xrightarrow{i \in I} f(x) = y \). Por lo tanto \( f(D) \) es \( \sigma \)-denso en \( Y \).
        \item Análogo al punto anterior restringiendo la función a \( f(X) \).
        \item Dado \( A \subseteq X \) y \( f \) continua, sea \( {(x_i)}_{i \in I} \) una red en \( A \). Si \( x \in A \) tal que \( x_i \xrightarrow{i \in I} x \) entonces \( f(x_i) \xrightarrow{i \in I} f(x) \) por la continuidad de \( f \). Por lo tanto \( f|_A : A \to Y \) es continua.
    \end{enumerate}
\end{proof}

\clearpage

\begin{ejercicio}{2}
    Sean \( X \) un conjunto y \( \tau \), \( \sigma \) dos topologías en \( X \). Probar que son equivalentes las proposiciones:
    \begin{enumerate}
        \item Se cumple \( \cl{A}^{\tau} \subseteq \cl{A}^{\sigma} \) para todo \( A \in \partes(X) \).
        \item \( \tau \)-convergencia de redes en \( X \) implica \( \sigma \)-convergencia de redes en \( X \).
        \item Se cumple \( \sigma \subseteq \tau \).
    \end{enumerate}
\end{ejercicio}
\begin{proof}
    Demostraremos (c) \( \Rightarrow \) (b) \( \Rightarrow \) (a) \( \Rightarrow \) (c).

    Empecemos por (c) \( \Rightarrow \) (b). Dado \( \sigma \subseteq \tau \) y \( {(x_i)}_{i \in I} \) una red que converge a \( x \) en \( \tau \) entonces \( \forall U \in \mathcal{O}^a_{\tau}(x) \) vale que
    \( \mathbf{x} \evto{U} \) y como \( \sigma \) tiene menos abiertos que \( \tau \) entonces \( \forall V \in \mathcal{O}^a_{\sigma}(x) \) vale que \( V \in \mathcal{O}^a_{\tau}(x) \) y por lo tanto \( \mathbf{x} \evto{V} \)
    \( \therefore {(x_i)}_{i \in I} \) converge a \( x \) en \( \sigma \).

    Veamos que (b) \( \Rightarrow \) (a). Dado \( A \in \partes(X) \) y \( x \in \cl{A}^{\tau} \) entonces existe una red \( \mathbf{x} \) tal que \( \mathbf{x} \evto{A} \) y \( \mathbf{x} \xrightarrow{i \in I} x \) en \( \tau \). Por la hipótesis se deduce que \( \mathbf{x} \xrightarrow{i \in I} x \) en \( \sigma \) y por lo tanto \( x \in \cl{A}^{\sigma} \).

    Finalmente para ver que (a) \( \Rightarrow \) (c) tomemos \( A \in \sigma \) entonces \( A^c \) es cerrado en \( \sigma \), entonces:  \[
        A^c = \cl{A^c}^{\sigma}
    \]
    Además, por hipotesis, sabemos que \( A^c \subseteq \cl{A^c}^{\tau} \subseteq \cl{A^c}^{\sigma} = A^c \implies \cl{A^c}^{\tau} = A^c \implies A \in \tau \)
\end{proof}

\begin{ejercicio}{3}
    Considerar en \( Y \) un orden total y la topología del orden respectiva. Sean \( f, g : X \to Y \) continuas. Entonces:
    \begin{enumerate}
        \item \( \{ x \in X : f(x) \leq g(x) \} \) es cerrado en \( X \).
        \item La función \( h(x) = \min \{ f(x), g(x) \} \) es continua.
    \end{enumerate}
    Notar que en la topología del orden, los abiertos son de la forma \( (a, b) = \{ y \in Y : a < y < b \} \), \( (a, +\infty) = \{ y \in Y : a < y \} \) o \( (-\infty, b) = \{ y \in Y : y < b \} \).
\end{ejercicio}
\begin{proof}
    Dado \( Y \) un orden total y la topología del orden respectiva, sean \( f,g : X \to Y \) continuas. Definamos: \[
        A = \{ x \in X : f(x) \leq g(x) \}
    \]
    Sea \( {(x_i)}_{i \in I} \) una red que vive en \( A \) con \( x_i \xrightarrow{i \in I} x \in X \). Nos preguntamos si \( x \in A \). En efecto,
    por la continuidad de \( f \) y \( g \) se deduce que \( f(x_i) \xrightarrow{i \in I} f(x) \) y \( g(x_i) \xrightarrow{i \in I} g(x) \).
    Supongamos que \( f(x) > g(x) \), entonces como \( Y \) es un orden total con la topología del orden entonces se pueden elegir dos abiertos que respeten el orden, es decir \( f(x) \in U = (c, +\infty) \), \( g(x) \in V = (-\infty, c) \) con \( c \) un punto intermedio \( f(x) > c > g(x) \), si no existiese tomo \( U = (g(x), +\infty) \), \( V = (-\infty, f(x)) \) y \( U > V \).
    Luego, \( f(x_i) \evto{U} \) y \( g(x_i) \evto{V} \) entonces por como elegimos \( U \) y \( V \) se deduce que \( f(x_i) > g(x_i) \) eventualmente, lo que contradice el hecho de que \( x_i \in A \). Por lo tanto \( f(x) \leq g(x) \) y \( A \) es cerrado.

    Para la función \( h \) basta notar que \[
        h(x) = \begin{cases}
            f(x) & \text{si } x \in \{ x \in X : f(x) \leq g(x) \} \\
            g(x) & \text{si } x \in \{ x \in X : g(x) \leq f(x) \}
        \end{cases} \]
    Como ambos conjuntos son cerrados y \( f \) y \( g \) son continuas y coinciden en la intersección entonces \( h \) es continua por el \emph{Lema del pegoteo}.
\end{proof}

\clearpage

\section*{Topologías}

\begin{ejercicio}{1}
    Sea \( P = \prod_{\alpha \in \mathbb{A}} (X_\alpha, \tau_\alpha) \) dotado de la topología producto \( \tau_P \). Para cada \( \alpha \in \mathbb{A} \) sea \( D_\alpha \subseteq X_\alpha \) denso.
    \begin{enumerate}
        \item Probar que el producto \( \prod_{\alpha \in \mathbb{A}} D_\alpha \) es denso en \( P \).
        \item Asumiendo que \( P \neq \varnothing \) se fija \( \mathbf{x} = {\{x_\alpha \}}_{\alpha \in \mathbb{A}} \in \prod_{\alpha \in \mathbb{A}} D_\alpha \) y se definen los conjuntos
              \[
                  D_F = \prod_{\alpha \in F} D_\alpha \times \prod_{\alpha \in \mathbb{A} - F} \{x_\alpha \} \subseteq P \qquad \text{para cada } F \in \pf(\mathbb{A})
              \]
              Mostrar que \( D = \bigcup_{F \in \pf(\mathbb{A})} D_F \) es denso en \( P \) y es más pequeño que el conjunto denso anterior.
    \end{enumerate}
\end{ejercicio}
\begin{proof}
    Sabemos que si \( U \in \tau_P \) entonces \( U = \prod_{\alpha \in F} U_\alpha \times \prod_{\alpha \in F^c} X_\alpha \) con \( F \in \pf{(\mathbb{A})} \) y \( U_\alpha \in \tau_\alpha \) para cada \( \alpha \in F \).
    Luego, como \( D_\alpha \) es denso en \( X_\alpha \) entonces \( \exists d_\alpha \in D_\alpha \cap U_\alpha \) para cada \( \alpha \in F \). Además, tomando un \( d_\alpha \in D_\alpha \cap X_\alpha = D_\alpha \) para cada \( \alpha \in F^c \) se deduce que
    \( {\{ d_\alpha \}}_{\alpha \in \mathbb{A}} \in U \implies \prod_{\alpha \in \mathbb{A}} D_\alpha \) es denso en \( P \).

    Para el segundo punto, siguiendo un razonamiento análogo se puede ver que para \( U \in \tau_P \) tomando el mismo \( F \) que antes, elijo \( D_F \) y armo el mismo \( {\{ d_\alpha \}}_{\alpha \in \mathbb{A}} \in U \) con \( {\{ d_\alpha \}}_{\alpha \in \mathbb{A}} \in D_F \subseteq D \). Por lo tanto, \( \bigcup_{F \in \pf{(\mathbb{A})}} D_F \) es denso en \( P \).
    Veamos ahora que es más chico que el conjunto anterior. Dado \( d \in D \) es claro que existe \( F \in \pf{(\mathbb{A})} \) tal que \( d \in D_F \). Entonces \( d_\alpha = x_\alpha \in D_\alpha \) si \( \alpha \notin F \) y \( d_\alpha \in D_\alpha \) si \( \alpha \in F \), es claro entonces que \( d \in \prod_{\alpha \in \mathbb{A}} D_\alpha \).
    Por lo tanto, \( D \subseteq \prod_{\alpha \in \mathbb{A}} D_\alpha \).
\end{proof}

\begin{ejercicio}{2}
    Sea \( P = \prod_{n \in \N} (X_n, \tau_n) \) dotado de la topología producto \( \tau_P \).
    \begin{enumerate}
        \item Si suponemos que todos los \( X_n \) son de tipo \( N_2 \) entonces \( P \) también lo es.
        \item Ídem con \( N_1 \) y separable.
        \item Muestre que no pasa lo mismo con Lindelöff.
    \end{enumerate}
\end{ejercicio}
\begin{proof}
    Veamos que si todos los \( X_n \) son \( N_2 \) entonces \( P \) también lo es. Dado \( \mathbf{x} = {(x_n)}_{n \in \N} \in P \) y \( U \in \mathcal{O}^a(\mathbf{x}) \) entonces \( U = \prod_{n \in F} U_n \times \prod_{n \in F^c} X_n \) con \( F \in \pf{(\N)} \) y \( U_n \in \mathcal{O}^a_{\tau_n}(x_n) \) para cada \( n \in F \).
    Luego, como \( X_n \) es \( N_2 \) entonces \( \exists V_n \in \mathcal{O}^a_{\tau_n}(x_n) \) tal que \( V_n \subseteq U_n \) para cada \( n \in F \). Entonces, \( V = \prod_{n \in F} V_n \times \prod_{n \in F^c} X_n \in \mathcal{O}^a(\mathbf{x}) \) y \( V \subseteq U \).
    Por lo tanto, \( P \) es \( N_2 \) pues construyo la base numerable \[
        \mathcal{B} = \left \{ \prod_{n \in F} V_n \times \prod_{n \in F^c} X_n : F \in \pf{(\N)}, V_n \in \mathcal{O}^a_{\tau_n}(x_n) \text{ con } V_n \subseteq U_n \text{ para cada } n \in F, U_n \in \mathcal{O}^a_{\tau_n}(x_n) \right \}
    \]
    pues \( \pf{(\N)} \) es numerable y cada \( \mathcal{O}^a_{\tau_n}(x_n) \) tiene una base numerable.

    Para el segundo punto, si cada \( X_n \) es \( N_1 \) entonces cada para cada \( {(x_n)}_{n \in \N} \in P \), la coordenada \( x_n \in X_n \) tiene una base numerable de entornos \( \mathcal{B}_n \).
    Luego, \[
        \mathcal{B} = \left \{ \prod_{n \in F} U_n \times \prod_{n \in F^c} X_n : F \in \pf{(\N)}, U_n \in \mathcal{B}_n \text{ para cada } n \in F \right \}
    \] es una base numerable de entornos de \( {(x_n)}_{n \in \N} \) y por lo tanto \( P \) es \( N_1 \). Por otro lado, si cada uno de los \( X_n \) es separable, significa que \( \exists D_n \subseteq X_n \)
    tal que \( \cl{D_n} = X_n \quad \forall n \in \N \). Luego, como \[
        \cl{\prod_{n \in \N} D_n} = \prod_{n \in \N} \cl{D_n} = \prod_{n \in \N} X_n = P
    \] entonces \( P \) es separable tomando \( D = \prod_{n \in \N} D_n \).

    Para ver que no vale con Lindelöff ver \emph{La máquina de generar contraejemplos} del apunte de Demetrio.
\end{proof}

\clearpage

\begin{ejercicio}{3}
    Sea \( f : \R \to \R^{\N} \) definida por \( f(t) = (t, t, t, \ldots) \). ¿Es continua si se considera en \( \R \) la topología usual y en \( \R^{\N} \) la topología caja?
\end{ejercicio}
\begin{proof}
    Consideremos \( \tau_d \) la topología usual de \( \R \) y \( \tau_b \) la topología de la caja en \( \R^{\N} \).
    Veamos primero que \( f \) no es continua en \( \tau_b \). Para ello, consideremos el abierto \( U = \prod_{n \in \N} (t - 1/n, t + 1/n) \) con \( t \in \R \). Entonces, \( f^{-1}(U) = \{t\} \) que no es un abierto de \( \R \). Por lo tanto, \( f \) no es continua.
    Sin embargo, veamos que \( f \) es continua en \( \tau_d \). Para ello, dada una sucesión \( {(t_i)}_{i \in \N} \) que converge a \( t \) en \( \tau_d \) entonces \( t_i \to t \) en la topología usual de \( \R \). Por lo tanto, \( f(t_i) = (t_i, t_i, t_i, \ldots) \to (t, t, t, \ldots) = f(t) \) en la topología de la caja pues cada coordenada converge a \( t \). Por lo tanto, \( f \) es continua.
\end{proof}

\begin{ejercicio}{4}
    Sea \( L \) el subconjunto de \( \R^{\N} \) formado por los \( (x_1, x_2, \ldots) \) con una cantidad finita de componentes no nulas. ¿Cuál es la clausura de \( L \) en \( \R^{\N} \) con la topología producto y con la topología caja?
\end{ejercicio}
\begin{proof}
    Sea \[ L = \left \{ (x_1, x_2, \ldots) \in \R^\N : \exists n_1, \ldots, n_k \in \N \, | \, x_{n_i} \neq 0 \text { para } i = 1, \ldots, k \text{ y } x_n = 0 \text{ para } n \notin \{n_1, \ldots, n_k\} \right \} \]
    Queremos encontrar \( \cl{L} \) con la topología producto y con la topología caja.
    En primer lugar, consideremos \( \tau_P \) la topología producto en \( \R^\N \). Podemos ver que \( \cl{L} = \R^\N \) tomando un abierto \( U = \prod_{\alpha \in F} U_\alpha \times \prod_{\alpha \notin F} X_\alpha \).
    Luego, tomamos el elemento de \( L \) que tiene en las finitas coordenadas \( \alpha \in F \) valores de \( U_\alpha \) y \( 0 \) en el resto de coordenadas.
    Ese elemento de \( L \) pertenece a \( U \) y por lo tanto \( U \cap L \neq \varnothing \quad \forall U \in \tau_P \therefore \cl{L} = \R^\N \) con la topología producto.

    En la topología caja tenemos que \( \cl{L} = L \). En efecto, veamos que \( X \setminus L \) es abierto. Sea \( (x_1, x_2, \ldots) \in X \setminus L \).
    Tenemos que existen solo finitas coordenadas donde \( x_n \neq 0 \). Si \( x_i = 0 \implies U_i = \R \), si \( x_i \neq 0 \implies U_i = (x_i - \varepsilon, x_i + \varepsilon) \) con \( \varepsilon > 0 \) tal que el entorno no contenga al cero.
    Luego, es fácil ver que \( x \in U = \prod_{i \in \N} U_i \subseteq X \setminus L \) por lo que todo punto es interior y \( X \setminus L \) es abierto \( \therefore L \) es cerrado.
\end{proof}

\begin{ejercicio}{5}
    Sea \( \Gamma \) un conjunto no-numerable. Se define
    \[
        \varphi(\Gamma) = \{ f : \Gamma \to \R : f(\lambda) = 0 \text{ para todo } \lambda \in \Gamma \text{ salvo para un conjunto } A \subseteq \Gamma \text{ numerable} \}
    \]
    Considerando que \( \varphi(\Gamma) \subseteq \R^{\Gamma} \) y en \( \R^{\Gamma} \) está la topología producto. Demostrar que \( \cl{\varphi(\Gamma)} = \R^{\Gamma} \).
\end{ejercicio}
\begin{proof}
    Dado \( \Gamma \) un conjunto no numerable, queremos ver que \( \cl{\varphi(\Gamma)} = \R^\Gamma \). Ello equivale a ver que el conjunto es denso, por lo que tomemos
    un abierto \( U = \prod_{\lambda \in F} U_\lambda \times \prod_{\lambda \notin F} \R \) con \( F \in \pf{(\Gamma)} \).
    Luego, tomamos la función \( f : \Gamma \to \R \) dada por \( f(\lambda) = 0 \) si \( \lambda \notin F \) y \( f(\lambda) \in U_\lambda \) si \( \lambda \in F \). Entonces, \( f \in U \cap \varphi(\Gamma) \) y por lo tanto \( U \cap \varphi(\Gamma) \neq \varnothing \quad \forall U \in \tau_P \therefore \cl{\varphi(\Gamma)} = \R^\Gamma \).
\end{proof}

\begin{ejercicio}{6}
    Sea \( (X, \tau) \) un espacio topológico y \( g : X \to Y \) suryectiva. ¿Cuál es la topología más pequeña en \( Y \) que hace que \( g \) sea continua?
\end{ejercicio}
\begin{proof}
    La topología más pequeña en \( Y \) que hace que \( g \) sea continua es \( \tau_{\text{ind}} = \{ \varnothing, Y \} \).
\end{proof}

\begin{ejercicio}{7}
    Sea \( g : \R \to S^1 \) dada por \( g(t) = e^{i2\pi t} \).
    \begin{enumerate}
        \item Como \( (\R, +) \) y \( (S^1, \cdot) \) son grupos, mostrar que \( g \) es un morfismo.
        \item Mostrar que \( \ker(g) = \Z \). Y que la relación de equivalencia \( \equiv \) dada por \( g \) es la congruencia módulo \( \Z \). Las clases de equivalencia son los conjuntos \( t + \Z \) para \( t \in [0, 1) \).
        \item Mostrar que \( \R/\Z \) es homeomorfo a \( S^1 \). (\( \R/\Z \) con la topología cociente y \( S^1 \) con la topología inducida por \( \C \)).
        \item Mostrar que \( g \) no es cerrada.
    \end{enumerate}
\end{ejercicio}
\begin{proof}
    Dada \( g : \R \to S^1 \) dada por \( g(t) = e^{i2\pi t} \).
    \begin{enumerate}
        \item Dado \( t, s \in \R \) entonces \( g(t+s) = e^{i2\pi (t+s)} = e^{i 2 \pi t} e^{i 2 \pi s} = g(t) g(s) \) por lo que \( g \) es un morfismo.
        \item Veamos que \( \ker(g) = \Z \). Dado \( t \in \R \) entonces \( g(t) = 1 \iff e^{i 2 \pi t} = 1 \iff i 2 \pi t = 2 \pi k i \) con \( k \in \Z \), por lo que \( t = k \in \Z \). Por lo tanto, \( \ker(g) = \Z \). Además, la relación de equivalencia dada por \( g \) es la congruencia módulo \( \Z \) pues \( t \equiv s \mod \Z \iff t - s \in \Z \iff g(t) = g(s) \).
        \item TODO
        \item Consideremos \[ A = \{ 1/n + n : n \in \N \} \subseteq \R \] que es cerrado por ser un conjunto formado por puntos aislados con distancia creciente entre ellos. Sin embargo, \[
                  g(A) = \{ e^{i 2 \pi (1/n + n)} = e^{i 2 \pi / n} : n \geq 2 \}
              \] que no es cerrado pues \( e^{i 2 \pi / n } \to 1 \) pero \( 1 \notin g(A) \).
    \end{enumerate}
\end{proof}

\begin{ejercicio}{8}
    Mostrar que existen funciones abiertas que no son cerradas y funciones cerradas que no son abiertas.
\end{ejercicio}
\begin{proof}
    La \( g \) del punto anterior es abierta, pero no es cerrada. Para una ejemplo de una función cerrada que no es abierta, consideremos \( f : \R \to \R \) dada por \[
        f(x) = \begin{cases}
            0 & x \leq 0 \\
            1 & x > 0
        \end{cases}
    \]
    Claramente no es abierta pues \( f((-\varepsilon, \varepsilon)) = \{0, 1\} \) que no es un abierto de \( \R \). Sin embargo, es cerrada pues para todo cerrado \( F \subseteq \R \) \( f(F) = \{ 0 \} \) ó \( f(F) = \{ 1 \} \) ó bien \( f(F) = \{ 0, 1 \} \) que son cerrados en \( \R \).
\end{proof}

\begin{ejercicio}{9}
    Sean \( Y = \{(x, y) : x = 0 \text{ ó } y = 0\} \) y \( g : \R^2 \to Y \) definida por
    \[
        g(x, y) = \begin{cases} (x, 0) & x \neq 0 \\ (0, y) & x = 0 \end{cases}
    \]
    Mostrar que \( g \) es suryectiva y que la topología cociente asociada \( \tau_g \) no es Hausdorff (\( \R^2 \) con la topología usual).
\end{ejercicio}
\begin{proof}
    Es claro que \( g \) es suryectiva pues dado \( (x, y) \in Y \) entonces \( g(x, 0) = (x, 0) \) si \( x \neq 0 \) y \( g(0, y) = (0, y) \) si \( x = 0 \).
    Para ver que \[ \tau_g = \{ U \subseteq Y : g^{-1}(U) \in \tau \} \] no es Hausdorff, tomemos \( (0, 1), (0, 2) \in Y \)
    y supongamos que existen \( U, V \in \tau_g \) con \( (0, 1) \in U \), \( (0, 2) \in V \) y \( U \cap V = \varnothing \). Entonces,
    \( g^{-1}(U) \) es un abierto de \( \R^2 \) que contiene a \( (0, 1) \) y \( g^{-1}(V) \) es un abierto de \( \R^2 \) que contiene a \( (0, 2) \). Como estamos en la topología usual,
    podemos decir que \[
        \exists \varepsilon > 0 : B((0, 1), \varepsilon) \subseteq g^{-1}(U) \quad \text{y} \quad \exists \delta > 0 : B((0, 2), \delta) \subseteq g^{-1}(V)
    \]
    Tomando \( 0 < x_0 < \min \{\varepsilon, \delta \} \) entonces el punto \( (x_0, 1) \in B((0, 1), \varepsilon) \) y \( (x_0, 2) \in B((0, 2), \delta) \).
    Sin embargo, \( g(x_0, 1) = (x_0, 0) = g(x_0, 2) \) lo que implica que \( (x_0, 1) \in g^{-1}(U) \cap g^{-1}(V) \) y por lo tanto \( U \cap V \neq \varnothing \). Por lo tanto, \( \tau_g \) no es Hausdorff.

\end{proof}

\clearpage

\section*{Métricas}

\begin{ejercicio}{1}
    Mostrar que las siguientes métricas en \( \R \) son topológicamente equivalentes pero no son equivalentes (a secas). Mostrar que \( (\R, d_2) \) no es completo.
    \[
        d_1(x, y) = |x - y| \qquad \text{y} \qquad d_2(x, y) = |f(x) - f(y)| \quad \text{siendo } f(t) = \frac{t}{1 + |t|}
    \]
\end{ejercicio}
\begin{proof}
    Veamos que no son equivalentes (a secas) i.e que no existen \[ m, M \in \R \quad : \quad \forall x, y \in \R \quad m d_2(x, y) \leq d_1(x, y) \leq M d_2(x, y) \] Para ello, consideremos
    \( d_1(x, x+1) = 1 \) y \( d_2(x, x+1) = \dfrac{1}{(1+x)(2+x)} \) si \( x > 0 \). Luego, es fácil ver que \( d_2(x, x+1) \to 0 \) cuando \( x \to +\infty \) por lo que ningún \( M > 0 \) puede cumplir el lado derecho de la ecuación.

    Para ver que son topológicamente equivalentes, consideremos la sucesión \( {(x_n)}_{n \in \N} \subseteq \R \) y \( x \in \R \), luego si \[
        d_1(x_n, x) = | x_n - x | \to 0
    \]
    \begin{align*}
        d_2(x_n, x) & = \left| \frac{x_n}{1+ |x_n|} - \frac{x}{1+|x|} \right | = \left| \frac{x_n + x_n | x | - (x + x | x_n |)}{(1 + |x_n|)(1 + |x|)} \right| \\
                    & \leq \left| (x_n - x) + x_n |x| - x |x_n| \right |                                                                                       \\
                    & \leq \left| (x_n - x) + x_n |x| - x_n |x_n| + x_n |x_n| - x |x_n| \right|                                                                \\
                    & \leq \left| x_n - x \right| + \left| x_n \right| \cdot \left| |x| - |x_n| \right| + \left| x_n - x \right| \cdot |x_n|                   \\
                    & \leq \left| x_n - x \right| + \left| x_n \right| \cdot \left| x - x_n \right| + \left| x_n - x \right| \cdot |x_n|                       \\
                    & \leq d(x_n, x) + 2 |x_n| \cdot d(x_n, x)                                                                                                 \\
                    & \leq d(x_n, x) (1 + 2 |x_n|)                                                                                                             \\
                    & \leq d(x_n, x) (1 + 2 K) \quad \text{ pues \( x_n \to x \) entonces \( |x_n| \) es acotada por alguna \( K > 0 \) }
    \end{align*}
    que tiende a \( 0 \) pues \( d_1(x_n, x) \to 0 \). Ahora, si \( d_2(x_n, x) \to 0 \) entonces \( \dfrac{x_n}{1 + |x_n|} \to \dfrac{x}{1 + |x|} \). Si \( x = 0 \) entonces \( x_n \to 0 \) y por lo tanto \( d_1(x_n, x) \to 0 \). Si \( x \neq 0 \) entonces \( x_n \to x \) y por lo tanto \( d_1(x_n, x) \to 0 \). Por lo tanto, son topológicamente equivalentes.

    Para ver que \( (\R, d_2) \) no es completo, consideremos la sucesión \( {(x_n)}_{n \in \N} \subseteq \R \) dada por \( x_n = n \). Entonces, \( d_2(x_n, x_m) = \dfrac{|n - m|}{(1 + n)(1 + m)} \to 0 \) cuando \( n, m \to +\infty \) por lo que \( {(x_n)}_{n \in \N} \) es de Cauchy. Sin embargo, no converge a ningún punto de \( \R \) pues \( d_2(x_n, x) = \left| \dfrac{n}{1 + n} - \dfrac{x}{1 + |x|} \right| \to 1 - \dfrac{x}{1 + |x|} > 0 \) para todo \( x \in \R \). Por lo tanto, \( (\R, d_2) \) no es completo.
\end{proof}

\begin{ejercicio}{2}
    Dado un conjunto \( X \) se define \( d(x, y) = \begin{cases} 1 & x \neq y \\ 0 & x = y \end{cases} \). Mostrar que \( d \) es una métrica. ¿Cuál es la topología \( \tau_d \)?
\end{ejercicio}
\begin{proof}
    Es claro que \( d \geq 0 \), \( d(x, y) = 0 \iff x = y \) por definición, \( d(x, y) = d(y, x) \) por simetría de la definición, y \( d(x, z) \leq d(x, y) + d(y, z) \) pues \( d(x, z) \) sólo puede ser \( 0 \) ó \( 1 \) y \( d(x, y) + d(y, z) \) sólo puede ser \( 0 \), \( 1 \) ó \( 2 \) y siempre se cumple la desigualdad triangular. Por lo tanto, \( d \) es una métrica.

    Para ver cómo es la topología \( \tau_d \) nos preguntamos que forma tienen las bolas de radio \( r \) con centro en \( x \). Entonces, dado \(  0 < r \leq 1 \) y \( x \in X \) entonces \[
        B(x, r) = \{ y \in X : d(x, y) < r \} = \{ x \}
    \]
    y si \( r > 1 \) entonces \[
        B(x, r) = \{ y \in X : d(x, y) < r \} = X
    \]
    Por lo tanto, la topología \( \tau_d \) es la topología discreta.
\end{proof}

\begin{ejercicio}{3}
    Mostrar que si \( (X, \tau) \) y \( (Y, \sigma) \) son homeomorfos y \( \tau \) es metrizable entonces \( \sigma \) también.
\end{ejercicio}
\begin{proof}
    Tenemos por hipótesis que existe un homeomorfismo \( f : X \to Y \) y una métrica \( d \) en \( X \) tal que \( \tau = \tau_d \).
    Entonces, definamos \( d' : Y \times Y \to \R \) dada por \( d'(y_1, y_2) = d(f^{-1}(y_1), f^{-1}(y_2)) \). Es fácil ver que \( d' \) es una métrica pues \( d' \geq 0 \), \( d'(y_1, y_2) = 0 \iff f^{-1}(y_1) = f^{-1}(y_2) \iff y_1 = y_2 \), \( d'(y_1, y_2) = d(f^{-1}(y_1), f^{-1}(y_2)) = d(f^{-1}(y_2), f^{-1}(y_1)) = d'(y_2, y_1) \), y \( d'(y_1, y_3) = d(f^{-1}(y_1), f^{-1}(y_3)) \leq d(f^{-1}(y_1), f^{-1}(y_2)) + d(f^{-1}(y_2), f^{-1}(y_3)) = d'(y_1, y_2) + d'(y_2, y_3) \).

    Finalmente, veamos que \( \tau_{d'} = \sigma \). Dado \( y \in Y \) y \( r > 0 \) entonces \( B_{d'}(y, r) = \{ z \in Y : d'(y, z) < r \} = \{ z \in Y : d(f^{-1}(y), f^{-1}(z)) < r \} = f(B_d(f^{-1}(y), r)) \) que es un abierto de \( Y \) pues \( f \) es un homeomorfismo. Por lo tanto, \( B_{d'}(y, r) \in \sigma \) y por lo tanto \( \tau_{d'} \subseteq \sigma \). Por otro lado, dado \( U \in \sigma \) entonces \( f^{-1}(U) \in \tau_d \) por ser \( f \) un homeomorfismo. Luego, dado \( x \in f^{-1}(U) \) y \( r > 0 \) tal que \( B_d(x, r) \subseteq f^{-1}(U) \) entonces \( f(B_d(x, r)) = B_{d'}(f(x), r) \subseteq U \) por lo que \( U \in \tau_{d'} \) y por lo tanto \( \sigma \subseteq \tau_{d'} \). Por lo tanto, \( \tau_{d'} = \sigma \) y \( \sigma \) es metrizable.
\end{proof}

\begin{ejercicio}{4}
    Sea \( P = \prod_{n \in \N} (X_n, \tau_n) \) dotado de la topología producto \( \tau_P \). Si cada \( \tau_n \) proviene de una métrica \( d_n \) en \( X_n \) entonces existe una métrica \( d_P \) en \( P \) tal que \( \tau_P = \tau_{d_P} \).
\end{ejercicio}
\begin{proof}
    Cabe hacer la aclaración de que si la métrica \( d_n \) en \( X_n \) genera a \( \tau_n \) entonces la métrica \( d_n' = \min \{ d_n, 1 \} \) es una métrica que también genera a \( \tau_n \).
    Entonces, dado \( P \) como arriba, podemos definir \( d_P : P \times P \to \R \) dada por \[
        d_P({(x_n)}_{n \in \N}, {(y_n)}_{n \in \N}) = \sum_{n \geq 1} \frac{d_n'(x_n, y_n)}{2^n}
    \]
    Veamos que \( d_P \) es una métrica. Es claro que \( d_P \geq 0 \) y \( d_P({(x_n)}, {(y_n)}) = 0 \iff d_n'(x_n, y_n) = 0 \quad \forall n \in \N \iff x_n = y_n \quad \forall n \in \N \iff {(x_n)} = {(y_n)} \). Además, \( d_P({(x_n)}, {(y_n)}) = d_P({(y_n)}, {(x_n)}) \) por simetría de cada \( d_n' \). Finalmente, dado \( {(x_n)}, {(y_n)}, {(z_n)} \in P \) entonces \( d_P({(x_n)}, {(z_n)}) = \sum_{n \geq 1} \frac{d_n'(x_n, z_n)}{2^n} \leq \sum_{n \geq 1} \frac{d_n'(x_n, y_n) + d_n'(y_n, z_n)}{2^n} = d_P({(x_n)}, {(y_n)}) + d_P({(y_n)}, {(z_n)}) \). Por lo tanto, \( d_P \) es una métrica.

    Demostremos ahora la inclusión de topologías \( \tau_P \subseteq \tau_{d_P} \). Sea \( U = \prod_{n \in F} U_n \times \prod_{n \notin F} X_n \) un abierto básico de la topología producto, donde \( F \subset \mathbb{N} \) es un conjunto finito y \( U_n \in \tau_n \) para cada \( n \in F \). Sea \( {(x_n)} \in U \), lo que implica que \( x_n \in U_n \) para todo \( n \in F \). Dado que cada \( \tau_n \) está generada por \( d_n' \), para cada \( n \in F \) existe un \( r_n > 0 \) tal que \( B_{d_n'}(x_n, r_n) \subseteq U_n \). Tomando \( r = \min_{n \in F} \frac{r_n}{2^n} \), que es estrictamente mayor que cero por ser \( F \) finito, afirmamos que \( B_{d_P}({(x_n)}, r) \subseteq U \). En efecto, si \( {(y_n)} \in B_{d_P}({(x_n)}, r) \), se tiene que:
    \[
        \sum_{n=1}^{\infty} \frac{d_n'(x_n, y_n)}{2^n} < r
    \]
    Luego, para cualquier \( k \in F \):
    \[
        \frac{d_k'(x_k, y_k)}{2^k} \leq \sum_{n=1}^{\infty} \frac{d_n'(x_n, y_n)}{2^n} < r \leq \frac{r_k}{2^k}
    \]
    De donde se deduce que \( d_k'(x_k, y_k) < r_k \), es decir, \( y_k \in B_{d_k'}(x_k, r_k) \subseteq U_k \). Como esto se cumple para todo \( k \in F \), concluimos que \( {(y_n)} \in U \). Así, \( B_{d_P}({(x_n)}, r) \subseteq U \), lo que prueba que \( U \in \tau_{d_P} \) y por lo tanto \( \tau_P \subseteq \tau_{d_P} \).

    Para ver la otra inclusión, \( \tau_{d_P} \subseteq \tau_P \), consideremos una bola abierta \( B_{d_P}({(x_n)}, r) \) con \( r > 0 \). Dado que la serie \( \sum_{n=1}^{\infty} \frac{1}{2^n} \) converge, existe un \( N \in \mathbb{N} \) suficientemente grande tal que:
    \[
        \sum_{n=N+1}^{\infty} \frac{1}{2^n} < \frac{r}{2}
    \]
    Definamos \( \delta = \frac{r}{2N} > 0 \) y consideremos el siguiente entorno abierto básico en la topología producto:
    \[
        V = B_{d_1'}(x_1, \delta) \times B_{d_2'}(x_2, \delta) \times \cdots \times B_{d_N'}(x_N, \delta) \times \prod_{n=N+1}^{\infty} X_n
    \]
    Es evidente que \( {(x_n)} \in V \in \tau_P \). Veamos ahora que \( V \subseteq B_{d_P}({(x_n)}, r) \). Si \( {(y_n)} \in V \), podemos acotar su distancia a \( {(x_n)} \) separando la serie en dos partes:
    \[
        d_P({(x_n)}, {(y_n)}) = \sum_{n=1}^{N} \frac{d_n'(x_n, y_n)}{2^n} + \sum_{n=N+1}^{\infty} \frac{d_n'(x_n, y_n)}{2^n}
    \]
    Para los primeros \( N \) términos, usamos que \( d_n'(x_n, y_n) < \delta \); para el resto, usamos el hecho de que \( d_n' \leq 1 \):
    \[
        d_P({(x_n)}, {(y_n)}) < \sum_{n=1}^{N} \frac{\delta}{2^n} + \sum_{n=N+1}^{\infty} \frac{1}{2^n} < \sum_{n=1}^{N} \delta + \frac{r}{2} = N\delta + \frac{r}{2} = N\left(\frac{r}{2N}\right) + \frac{r}{2} = r
    \]
    Por lo tanto, \( {(y_n)} \in B_{d_P}({(x_n)}, r) \), lo que demuestra que \( V \subseteq B_{d_P}({(x_n)}, r) \). Esto prueba que \( B_{d_P}({(x_n)}, r) \in \tau_P \), de donde \( \tau_{d_P} \subseteq \tau_P \). Finalmente, concluimos que \( \tau_P = \tau_{d_P} \).
\end{proof}

\begin{ejercicio}{5}
    Sea \( L \) el subconjunto de \( \R^{\N} \) formado por las sucesiones con una cantidad finita de componentes no nulas. ¿Cuál es la clausura en \( \R^{\N} \) respecto a la topología uniforme?
\end{ejercicio}
\begin{proof}
    Propongo que \( \cl{L} = c_0 = \{ \text{sucesiones convergentes a 0} \} \). En efecto,
    veamos que \( \cl{L} \subseteq c_0 \). Dados \( f \in \cl{L} \) y \( \varepsilon > 0 \) existe \( g \in L \) tal que \( d_\infty(f, g) < \varepsilon \). Dado que \( g \) tiene sólo finitas componentes no nulas,
    existe un \( n_0 \in \N \) tal que \( g(n) = 0 \quad \forall n \geq n_0 \). Luego, para todo \( n \geq n_0 \) se tiene que \( |f(n)| = |f(n) - g(n)| < \varepsilon \) por lo que \( f(n) \to 0 \) y por lo tanto \( f \in c_0 \).

    Por último, veamos que \( c_0 \subseteq \cl{L} \). Dado \( f \in c_0 \) y \( \varepsilon > 0 \) existe un \( n_0 \in \N \) tal que \( |f(n)| < \varepsilon \quad \forall n \geq n_0 \). Entonces, la función \( g : \N \to \R \) dada por \( g(n) = f(n) \) si \( n < n_0 \) y \( g(n) = 0 \) si \( n \geq n_0 \) es una función con sólo finitas componentes no nulas, es decir, \( g \in L \). Además, \( d_\infty(f, g) = \sup_{n} |f(n) - g(n)| = \sup_{n} |f(n)| < \varepsilon \). Por lo tanto, \( f \in B_\infty(g, \varepsilon) \cap L \neq \varnothing \), lo que implica que \( f \in \cl{L} \). Por lo tanto, \( c_0 = \cl{L} \).
\end{proof}

\begin{ejercicio}{6}
    Considerando las funciones \( E, F : C([0,1], \R) \to \R \) definidas por
    \[
        E(f) = f(0) \qquad F(f) = \int_0^1 f(t)\, dt
    \]
    \begin{enumerate}
        \item Demostrar que son continuas considerando en \( C([0,1], \R) \) la métrica uniforme.
        \item Demostrar que \( F \) es continua pero \( E \) no lo es si se considera en \( C([0,1], \R) \) la métrica
              \[
                  d_1(f, g) = \int_0^1 |f(t) - g(t)|\, dt
              \]
    \end{enumerate}
\end{ejercicio}
\begin{proof}
    Como estamos en espacios métricos para probar la continuidad de las funciones \( E, F \) podemos aplicar un argumento de \( \varepsilon \)-\( \delta \). Entonces, fijado \( f \in C([0, 1]) \) y \( \varepsilon > 0 \), tomando
    \( \delta = \varepsilon \) se tiene que si: \[ d_\infty(f, g) = \sup_{t \in [0, 1]} |f(t) - g(t)| < \delta = \varepsilon \] entonces \[
        |E(f) - E(g)| = |f(0) - g(0)| < d_\infty(f, g) < \varepsilon
    \]
    por lo tanto \( E \) es continua. Además, si \( d_\infty(f, g) < \varepsilon \) entonces \[
        |F(f) - F(g)| = \left| \int_0^1 f(t)\, \mathrm{d}t - \int_0^1 g(t)\, \mathrm{d}t \right| = \left| \int_0^1 (f(t) - g(t))\, \mathrm{d}t \right| \leq \int_0^1 |f(t) - g(t)|\, \mathrm{d}t < d_\infty(f, g) < \varepsilon
    \]
    por lo que \( F \) es continua.

    Si ahora utilizamos la métrica \( d_1 \) entonces, dado \( f \in C([0, 1]) \) y \( \varepsilon > 0 \), tomando \( \delta = \varepsilon \) se tiene que si: \[
        d_1(f, g) = \int_0^1 |f(t) - g(t)|\, dt < \delta = \varepsilon
    \]
    entonces \[
        |F(f) - F(g)| = \left| \int_0^1 f(t)\, \mathrm{d}t - \int_0^1 g(t)\, \mathrm{d}t \right| = \left| \int_0^1 (f(t) - g(t))\, \mathrm{d}t \right| \leq \int_0^1 |f(t) - g(t)|\, \mathrm{d}t = d_1(f, g) < \varepsilon
    \]
    por lo que \( F \) es continua. Sin embargo, \( E \) no es continua respecto de \( d_1 \). Para ver esto, fijada la \( f \in C([0, 1]) \) y \( \varepsilon > 0 \), para cualquier \( \delta > 0 \) se puede construir una función tal que
    \[
        |E(f) - E(g)| = |f(0) - g(0)| \geq \varepsilon
    \]
    sin embargo \( d_1(f, g) < \delta \). Esta función se construye tomando un punto inicial lo suficientemente lejos, pero luego tomando un pico lo suficientemente empinado tal que la integral de la diferencia entre \( f \) y \( g \) sea menor que \( \delta \).
    Para ello sólo necesitamos que la función \( g \) tenga un pico cuya diferencia con \( f \) sea de altura \( \varepsilon \) y base \( \frac{\delta}{\varepsilon} \) y luego se pegue a \( f \) en el resto del intervalo.
\end{proof}

\begin{ejercicio}{7}
    Mostrar que en un espacio métrico, la propiedad `separable' es hereditaria.
\end{ejercicio}
\begin{proof}
    Si \( X \) es un espacio métrico separable por la \textbf{Proposición 2.4.2} es es \( N_2 \), luego por la \textbf{Proposición 2.4.15} \( N_2 \) es hereditario y entonces \( Y \) es \( N_2 \). Pero \( Y \) es subespacio también métrico, entonces de nuevo por la \textbf{Proposición 2.4.2} es separable.
\end{proof}

\begin{ejercicio}{8}
    Mostrar que \( X \) es separable respecto a dos métricas \( d_1 \) y \( d_2 \) si y sólo si es separable respecto a la métrica \( d = \max \{ d_1, d_2 \} \).
\end{ejercicio}
\begin{proof}
    \( (\impliedby) \) Por hipótesis existe un \( A \subseteq X \) denso y numerable. Dado \( \varepsilon > 0 \) y \( x \in X \) entonces \( \exists a \in A \, : \, d(x, a) < \varepsilon \implies d_1(x, a) < \varepsilon \) y \( d_2(x, a) < \varepsilon \therefore A \) es denso numerable con respecto a \( d_1, d_2 \).

    \( (\implies) \) Por hipótesis \( X \) es un espacio métrico separable respecto a \( d_1 \) y a \( d_2 \). Entonces, \( X \) es \( N_2 \) con respecto a ambas métricas y esto nos dice que existen \( \mathcal{B}_1, \mathcal{B}_2  \) bases numerables de \( X \) con respecto a \( d_1 \) y \( d_2 \) respectivamente. Entonces, es fácil ver que \( \mathcal{B} = \{ B_1 \cap B_2 : B_1 \in \mathcal{B}_1, B_2 \in \mathcal{B}_2 \} \) es una base numerable de \( X \) con respecto a \( d \).
    Por lo tanto, \( X \) es \( N_2 \) y finalmente separable con respecto a \( d \).
\end{proof}

\clearpage

\section*{Funciones Continuas (Lema de Urysohn/Teorema de Tietze)}

\begin{ejercicio}{1}
    Sea \( (X, \tau) \) un espacio topológico. Probar que:
    \begin{enumerate}
        \item Si \( X \) es normal entonces \( X \) es completamente regular.
        \item La clase completamente regular es hereditaria.
        \item El producto de espacios completamente regulares sigue siendo completamente regular.
    \end{enumerate}
\end{ejercicio}
\begin{proof}
    Para ver que normal implica CR quiero ver que para todo par \( (x, U) \in X \times \tau \) con \( x \in U \) existe una función continua \( f : X \to [0, 1] \) tal que \( f(x) = 1 \) y \( f(y) = 0 \) para todo \( y \notin U \).
    Dado \( (x, U) \in X \times \tau \) con \( x \in U \), como \( X \) es normal, en particular \( T_1 \) entonces \( \{ x \} \) es cerrado y \( U^c \) también por ser complemento de un abierto.
    Luego, por el \emph{Lema de Urysohn} existe una función continua \( f : X \to [0, 1] \) tal que \( f(x) = 1 \) y \( f(y) = 0 \) para todo \( y \notin U \). Por lo tanto, \( X \) es completamente regular.

    Para ver que CR es hereditaria, consideremos \( Y \subseteq X \) un subespacio de \( X \) y \( (y, U) \in Y \times \tau_Y \) con \( y \in U \).
    Entonces, existe un \( V \in \tau \) tal que \( U = V \cap Y \). Dado que \( X \) es CR, existe una función continua \( f : X \to [0, 1] \) tal que \( f(y) = 1 \) y \( f(x) = 0 \) para todo \( x \notin V \).
    Entonces, la función \( f|_Y : Y \to [0, 1] \) dada por \( f|_Y(z) = f(z) \) es continua pues es la restricción de una función continua a un subespacio. Además, \( f|_Y(y) = f(y) = 1 \) y \( f|_Y(z) = f(z) = 0 \) para todo \( z \notin U \). Por lo tanto, \( Y \) es completamente regular.

    Sean \( X_1, X_2 \) completamente regulares, veamos que \( X_1 \times X_2 \) es completamente regular. En efecto, dado \( x = (x_1, x_2) \) y \( U = U_1 \times U_2 \) con \( x \in U \) entonces \( x_1 \in U_1 \) y \( x_2 \in U_2 \). Dado que \( X_1, X_2 \) son completamente regulares, existen funciones continuas \( f_1 : X_1 \to [0, 1] \) y \( f_2 : X_2 \to [0, 1] \) tal que \( f_1(x_1) = 1 \), \( f_1(y) = 0 \) para todo \( y \notin U_1 \), \( f_2(x_2) = 1 \), y \( f_2(z) = 0 \) para todo \( z \notin U_2 \). Entonces, la función \( f : X_1 \times X_2 \to [0, 1] \) dada por \( f(z, w) = f_1(z) f_2(w) \) es continua pues es el producto de funciones continuas. Además, \( f(x) = f(x_1, x_2) = f_1(x_1) f_2(x_2) = 1\cdot 1 = 1 \), y si \( (z, w) \notin U = U_1 \times U_2 \), entonces o bien \( z \notin U_1 \) o bien \( w \notin U_2 \). En el primer caso, se tiene que \( f(z, w) = f_1(z)f_2(w) = 0\cdot f_2(w) = 0\), y en el segundo caso se tiene que \( f(z, w) = f_1(z)f_2(w) = f_1(z)\cdot 0 = 0\). Por lo tanto, \( X_1\times X_2 \) es completamente regular.
\end{proof}

\begin{ejercicio}{2}
    Sea \( (X, \tau) \) un espacio topológico completamente regular.
    \begin{enumerate}
        \item Probar que dada una red \( \mathbf{x} = {\{x_i\}}_{i \in I} \) en \( X \) y \( x \in X \) se cumple
              \[
                  x_i \xrightarrow{i \in I} x \iff f(x_i) \xrightarrow{i \in I} f(x) \quad \text{para toda } f \in C(X, [0,1])
              \]
        \item Probar que la topología original \( \tau \) en \( X \) no es otra que la topología inicial asociada a \( C(X, [0,1]) \).
    \end{enumerate}
\end{ejercicio}
\begin{proof}
    La ida es clara, si \( x_i \to x \) entonces \( f(x_i) \to f(x) \) para toda \( f \in C(X, [0,1]) \) por continuidad de \( f \). Para la vuelta, supongamos que \( f(x_i) \to f(x) \) para toda \( f \in C(X, [0,1]) \) pero \( x_i \not\to x \). Entonces, existe un abierto \( U \) con \( x \in U \) tal que \( x_i \notin U \) para toda \( i \in I \). Dado que \( X \) es completamente regular, existe \( f \in C(X, [0, 1]) \) tal que \( f(x) = 1 \) y \( f(y) = 0 \) para todo \( y \notin U \). Entonces, \( f(x_i) = 0 \) para toda \( i \in I \) pues los \( x_i \notin U \), pero \( f(x) = 1 \), lo que contradice la hipótesis de que \( f(x_i) \to f(x) \). Por lo tanto, \( x_i \to x \).
\end{proof}

\begin{ejercicio}{3}
    Sea \( X \) un espacio topológico completamente regular. Considerando \( C(X, \R) \subseteq \R^X \) demostrar que \( C(X, \R) \) es denso respecto de la topología producto en \( \R^X \).
\end{ejercicio}
\begin{proof}
    Para ver que \( C(X, \R) \) es denso con respecto a la topología producto, dado \( f \in \R^X \) y un entorno abierto básico \[ U = \prod_{x \in F} U_x \times \prod_{x \notin F} \R \] con \( F \subseteq X \) finito y
    \( U_x \subseteq \R \) abierto para cada \( x \in F \), quiero encontrar una función \( g \in C(X, \R) \) tal que \( g \in U \). Dado que \( X \) es CR, para cada \( x \in F \) existe una función continua
    \( f_x : X \to [0, 1] \) tal que \( f_x(x) = 1 \) y \( f_x(y) = 0 \) para todo \( y \neq x \). Entonces, la función \( g : X \to \R \) dada por
    \[
        g(z) = \sum_{x \in F} f_x(z) f(x)
    \]
    es continua pues es una suma finita de funciones continuas. Además, si \( z \in F \) entonces \( g(z) = f(z) \) pues \( f_x(z) = 1 \) para \( x = z \) y \( f_x(z) = 0 \) para \( x \neq z \). Por lo tanto, \( g(z) \in U_z \) para cada \( z \in F \), lo que implica que \( g \in U \). Así, \( C(X, \R) \cap U \neq \varnothing \), lo que demuestra que \( C(X, \R) \) es denso en \( \R^X \).
\end{proof}

\begin{ejercicio}{4}
    Sea \( (X, \tau) \) un espacio topológico regular y sean \( \mathbf{x} = {(x_i)}_{i \in I} \) e \( \mathbf{y} = {(y_j)}_{j \in J} \) dos redes en \( X \) ambas convergentes. Entonces son equivalentes:
    \begin{itemize}
        \item Las dos redes tienen el mismo límite.
        \item La red \( \{\mathbf{x}, \mathbf{y}\} = {(\{x_i, y_j\})}_{(i,j) \in I \times J} \) cumple \( \{\mathbf{x}, \mathbf{y}\} \evto{U} \) para todo abierto \( U \subseteq X \times X \) tal que \( \Delta \subseteq U \). El orden en \( I \times J \) es en las dos entradas a la vez: \( (i,j) \leq (k,l) \iff i \leq_I k \) y \( j \leq_J l \).
    \end{itemize}
\end{ejercicio}
\begin{proof}
    Sea \( \Delta = \{ (x, x) : x \in X \} \) la diagonal de \( X \times X \). Dado \( (x, y) \in X \times X \), entonces \( (x, y) \in \Delta \) si y sólo si \( x = y \).

    \( (\implies) \) Supongamos que las dos redes tienen el mismo límite \( z \). Sea \( U \subseteq X \times X \) un abierto tal que \( \Delta \subseteq U \). Como \( (z, z) \in \Delta \subseteq U \), tenemos que \( (z, z) \in U \).
    Por definición de la topología producto, existen abiertos \( W_1, W_2 \subseteq X \) tales que \( (z, z) \in W_1 \times W_2 \subseteq U \). Sea \( W = W_1 \cap W_2 \). Entonces \( W \) es un entorno abierto de \( z \) y \( W \times W \subseteq U \).
    Dado que ambas redes convergen a \( z \), existen índices \( i_0 \in I \) y \( j_0 \in J \) tales que \( x_i \in W \) para toda \( i \geq i_0 \) y \( y_j \in W \) para toda \( j \geq j_0 \). Por lo tanto, si \( (i, j) \geq (i_0, j_0) \), se tiene que \( (x_i, y_j) \in W \times W \subseteq U \). Esto demuestra que \( \{\mathbf{x}, \mathbf{y}\} \evto{U} \).

    \( (\impliedby) \) Supongamos que \( \{\mathbf{x}, \mathbf{y}\} \evto{U} \) para todo abierto \( U \subseteq X \times X \) tal que \( \Delta \subseteq U \). Por hipótesis, ambas redes son convergentes. Sean \( x_i \to a \) y \( y_j \to b \). Queremos demostrar que \( a = b \).

    Supongamos por el absurdo que \( a \neq b \). Asumiendo que los puntos son cerrados (es decir, el espacio es \( T_1 \) y por ser regular resulta Hausdorff), existen entornos abiertos disjuntos \( V_a \) y \( V_b \) de \( a \) y \( b \) respectivamente.
    Dado que \( X \) es regular, existe un entorno abierto \( W_a \) de \( a \) tal que \( \cl{W_a} \subseteq V_a \), y un entorno abierto \( W_b \) de \( b \) tal que \( \cl{W_b} \subseteq V_b \). Como \( V_a \cap V_b = \emptyset \), entonces \( \cl{W_a} \cap \cl{W_b} = \emptyset \).

    Consideremos el conjunto \( U = (X \times X) \setminus (\cl{W_a} \times \cl{W_b}) \). Como \( \cl{W_a} \) y \( \cl{W_b} \) son cerrados, su producto cartesiano es cerrado en \( X \times X \), por lo que \( U \) es un conjunto abierto.
    Además, como \( \cl{W_a} \cap \cl{W_b} = \emptyset \), no existe ningún \( x \in X \) tal que \( (x, x) \in \cl{W_a} \times \cl{W_b} \). En consecuencia, \( \Delta \subseteq U \).

    Por hipótesis, la red \( \{\mathbf{x}, \mathbf{y}\} \) está eventualmente en \( U \). Sin embargo, como \( x_i \to a \) y \( y_j \to b \), eventualmente se tiene que \( x_i \in W_a \subseteq \cl{W_a} \) y \( y_j \in W_b \subseteq \cl{W_b} \). Es decir, eventualmente \( (x_i, y_j) \in \cl{W_a} \times \cl{W_b} \), lo que implica que la red está eventualmente fuera de \( U \).
    Esta contradicción proviene de suponer \( a \neq b \). Por lo tanto, \( a = b \), y las dos redes tienen el mismo límite.
\end{proof}

\begin{ejercicio}{5}
    Sea \( X \) un espacio regular con una base numerable. Sea \( U \) un abierto de \( X \).
    \begin{enumerate}
        \item Mostrar que \( U \) es igual a la unión numerable de conjuntos cerrados de \( X \).
        \item Mostrar que existe una \( f \in C(X, [0,1]) \) tal que \( f(x) > 0 \) para los \( x \in U \) y \( f(x) = 0 \) para los \( x \notin U \).
    \end{enumerate}
\end{ejercicio}
\begin{proof}
    Dado \( U \in \tau \), como \( X \) es regular, para cada \( x \in U \) existe un entorno abierto \( V_x \) de \( x \) tal que \( x \in V_x \subseteq \cl{V_x} \subseteq U \). Luego, como tengo una base numerable
    existe \( x \in B_{n_x} \subseteq V_x \) para algún \( n_x \in \N \). Entonces, \( x \in \cl{B_{n_x}} \subseteq \cl{V_x} \subseteq U \). Por lo tanto, \[
        U = \bigcup_{x \in U} \cl{B_{n_x}}
    \]
    y cada \( \cl{B_{n_x}} \) es un conjunto cerrado de \( X \). Además, como \( B_{n_x} \) pertenece a una base numerable, entonces \( n_x \) sólo puede tomar valores en un subconjunto numerable de \( \N \), lo que implica que la unión anterior es numerable.

    Para ver la segunda parte, dado \( U \in \tau \), por la primera parte existe una sucesión \( (F_n)_{n \in \N} \) de conjuntos cerrados de \( X \) tal que \( U = \bigcup_{n \in \N} F_n \). Entonces, la función \( f : X \to [0, 1] \) dada por
    \[
        f(x) = \sum_{n=1}^{\infty} \frac{\chi_{F_n}(x)}{2^n}
    \] es continua pues es la suma de una serie de funciones continuas (cada \( \chi_{F_n} \) es continua por ser la función característica de un conjunto cerrado). Además, si \( x \in U \) entonces \( x \in F_{n_0} \) para algún \( n_0 \in \N \), lo que implica que \( f(x) \geq \frac{1}{2^{n_0}} > 0 \). Por otro lado, si \( x \notin U \) entonces \( x \notin F_n \) para todo \( n \in \N \), lo que implica que \( f(x) = 0 \). Por lo tanto, \( f \) es la función continua que cumple las condiciones requeridas.
\end{proof}

\begin{ejercicio}{6}
    Demostrar el \emph{Teorema Fuerte de Urysohn}: Dados \( A \) y \( B \) conjuntos cerrados disjuntos de un espacio topológico normal. Entonces existe \( f \in C(X, [0,1]) \) tal que \( f^{-1}(0) = A \) y \( f^{-1}(1) = B \) si y sólo si \( A \) y \( B \) son \( G_\delta \) conjuntos en \( X \).

    \textit{Definición}: Un conjunto \( G_\delta \) de \( X \) es aquel que puede escribirse como intersección numerable de conjuntos abiertos.
\end{ejercicio}

\begin{proof}
    Supongamos primero que existe \( f \in C(X, [0, 1]) \) tal que \( f^{-1}(0) = A \) y \( f^{-1}(1) = B \), y veamos que \( A \) y \( B \) son \( G_\delta \) conjuntos. Para cada \( n \in \N \), definamos los conjuntos:
    \[
        U_n = f^{-1}\left( \left[ 0, \frac{1}{n} \right) \right) \quad \text{y} \quad V_n = f^{-1}\left( \left( 1 - \frac{1}{n}, 1 \right] \right)
    \]
    Como \( f \) es continua y los intervalos son abiertos en la topología subespacio de \( [0,1] \), \( U_n \) y \( V_n \) son abiertos en \( X \). Luego, \( A = \bigcap_{n=1}^{\infty} U_n \) y \( B = \bigcap_{n=1}^{\infty} V_n \), lo que implica que \( A \) y \( B \) son conjuntos \( G_\delta \).

    Veamos ahora la otra implicación. Supongamos que
    \[
        A = \bigcap_{n=1}^{\infty} U_n \quad \text{y} \quad B = \bigcap_{n=1}^{\infty} V_n \quad \text{con } U_n, V_n \in \tau \, (\forall n \in \N)
    \]
    y veamos que existe \( f \in C(X, [0, 1]) \) tal que \( f^{-1}(0) = A \) y \( f^{-1}(1) = B \).

    Primero, construimos una función que se anule exactamente en \( A \). Como \( A \subseteq U_n \), el conjunto \( X \setminus U_n \) es un cerrado disjunto de \( A \). Por el Lema de Urysohn (ya que \( X \) es normal), para cada \( n \in \N \) existe \( g_n \in C(X, [0, 1]) \) tal que \( g_n(A) = \{0\} \) y \( g_n(X \setminus U_n) = \{1\} \). Definimos la función \( g : X \to [0, 1] \) como:
    \[
        g(x) = \sum_{n=1}^{\infty} \frac{g_n(x)}{2^n}
    \]
    Esta función es continua por ser suma de una serie uniformemente convergente de funciones continuas (Weierstrass). Si \( x \in A \), entonces \( g_n(x) = 0 \) para todo \( n \in \N \), luego \( g(x) = 0 \). Si \( x \notin A \), existe \( n_0 \in \N \) tal que \( x \notin U_{n_0} \), por lo que \( g_{n_0}(x) = 1 \). En consecuencia, \( g(x) \geq \frac{1}{2^{n_0}} > 0 \). Por lo tanto, \( g^{-1}(0) = A \).

    De manera análoga, construimos una función \( h : X \to [0, 1] \) para el conjunto \( B \). Para cada \( n \in \N \), separamos \( B \) de \( X \setminus V_n \) obteniendo funciones \( h_n \), y definimos \( h(x) = \sum_{n=1}^{\infty} \frac{h_n(x)}{2^n} \), lo que nos garantiza que \( h \) es continua y \( h^{-1}(0) = B \).

    Finalmente, definimos \( f : X \to [0,1] \) como:
    \[
        f(x) = \frac{g(x)}{g(x) + h(x)}
    \]
    Como \( A \cap B = \emptyset \), no existe ningún \( x \in X \) donde \( g(x) = 0 \) y \( h(x) = 0 \) simultáneamente. Por ende, el denominador nunca se anula y \( f \) es continua.

    Evaluando los conjuntos de nivel:
    Si \( x \in A \), entonces \( g(x) = 0 \) y \( h(x) > 0 \), por lo que \( f(x) = 0 \). Si \( f(x) = 0 \), debe ser \( g(x) = 0 \), luego \( x \in A \). Así, \( f^{-1}(0) = A \).
    Si \( x \in B \), entonces \( h(x) = 0 \) y \( g(x) > 0 \), por lo que \( f(x) = \frac{g(x)}{g(x)} = 1 \). Si \( f(x) = 1 \), entonces \( g(x) = g(x) + h(x) \implies h(x) = 0 \implies x \in B \). Así, \( f^{-1}(1) = B \).
\end{proof}

\begin{ejercicio}{7}
    Mostrar que el \emph{Teorema de Tietze} implica el \emph{Lema de Urysohn}.
\end{ejercicio}

\begin{proof}
    Recordemos que el \emph{Lema de Urysohn} establece que si \( X \) es un espacio normal y \( A, B \subseteq X \) son conjuntos cerrados disjuntos, entonces existe una función continua \( f : X \to [0, 1] \) tal que \( f(A) = \{0\} \) y \( f(B) = \{1\} \).

    Por otro lado, el \emph{Teorema de Tietze} establece que si \( X \) es un espacio normal y \( C \subseteq X \) es un conjunto cerrado, entonces cualquier función continua y acotada \( g : C \to [-1, 1] \) se puede extender a una función continua \( f : X \to [-1, 1] \) tal que \( f|_C = g \) y \( \left\| f \right\|_\infty = \left\| g \right\|_\infty \). Es importante notar que esto aplica equivalentemente para cualquier intervalo cerrado, en particular \( [0, 1] \).

    Teniendo la hipótesis de que \( X \) es normal, tomemos \( A \) y \( B \) los conjuntos cerrados disjuntos del \emph{Lema de Urysohn} y construyamos la función \( f \) utilizando el \emph{Teorema de Tietze}.

    Como \( A \) y \( B \) son cerrados, su unión \( C = A \cup B \) también es un conjunto cerrado en \( X \). Definamos la función \( g : A \cup B \to [0, 1] \) dada por:
    \[
        g(x) = \begin{cases}
            0 & \text{si } x \in A \\
            1 & \text{si } x \in B
        \end{cases}
    \]
    Esta función es continua pues \( A \) y \( B \) son conjuntos cerrados disjuntos, por lo que podemos aplicar el \emph{Lema del pegoteo}. Claramente, \( g \) es acotada ya que su imagen está contenida en el intervalo \( [0, 1] \).

    Por el \emph{Teorema de Tietze}, existe una función continua \( f : X \to [0, 1] \) tal que \( f|_{A \cup B} = g \) y \( \left \| f \right \|_\infty = \left \| g \right \|_\infty \). En particular, por la definición de \( g \), tenemos que \( f(A) = \{0\} \) y \( f(B) = \{1\} \), lo que demuestra el \emph{Lema de Urysohn}.
\end{proof}

\begin{ejercicio}{8}
    Demostrar que un espacio topológico normal y conexo con más de un punto es no-numerable.
\end{ejercicio}
\begin{proof}
    Sea \( X \) un espacio topológico normal y conexo con más de un punto. Supongamos por el absurdo que \( X \) es numerable.
    Dado que \( X \) tiene más de dos puntos distintos, entonces puedo tomar \( x, y \in X \) con \( x \neq y \). Como \( X \) es normal, es \( T_1 \) y los singueletes son cerrados. Por lo tanto, aplicando el
    \emph{Lema de Urysohn} para los conjuntos cerrados disjuntos \( \{ x \} \) y \( \{ y \} \), existe una función continua \( f : X \to [0, 1] \) tal que \( f(x) = 0 \) y \( f(y) = 1 \). Como \( X \) es conexo,
    la imagen de \( f \) también lo es, lo que implica que \( f(X) \) es un intervalo conectado que contiene a \( 0 \) y a \( 1 \), pero esto solo puede ocurrir si \( f(X) = [0, 1] \). Sin embargo, \( X \) es numerable,
    lo que implica que \( f(X) \) es numerable, lo que contradice el hecho de que \( f(X) = [0, 1] \) es no numerable. Por lo tanto, \( X \) no puede ser numerable.
\end{proof}

\begin{ejercicio}{9}
    Mostrar que si \( Y \) es normal y \( \mathcal{B} \) es una base de su topología. Entonces existe un embedding de \( Y \) en \( {[0,1]}^J \) siendo \( J \) algún subconjunto de \( \mathcal{B} \times \mathcal{B} \).
\end{ejercicio}
\begin{proof}
    Definimos el conjunto de índices \( J \) como:
    \[
        J := \{ (U, V) \in \mathcal{B} \times \mathcal{B} : \cl{U} \subseteq V \}
    \]
    Para cada par \( (U,V) \in J \), los conjuntos \( \cl{U} \) y \( Y \setminus V \) son cerrados y disjuntos en Y. Por el \emph{Lema de Urysohn}, existe una función continua \( f_{U,V} : Y \to [0,1] \) tal que \( f_{U,V}(\cl{U}) = \{0\} \) y \( f_{U,V}(Y \setminus V) = \{1\} \).

    Definimos la función \( F : Y \to {[0, 1]}^J \) dada por \( F(y) = {\left( f_{U, V}(y) \right)}_{(U, V) \in J} \).
    Dado que cada función coordenada \( f_{U, V} \) se sigue que \( F \) es continua.

    1. Inyectividad de \( F \):
    Sean \( y, z \in Y \), \( y \neq z \). Como \( Y \) es \( T_1 \), el conjunto \( Y \setminus \{ z \} \) es un abierto que contiene a \( y \).
    Por ser \( Y \) un espacio regular, existe un abierto \( W \) tal que \( y \in W \subseteq \cl{W} \subseteq Y \setminus \{z\} \).

    Como \( \mathcal{B} \) es base, existe un elemento \( V \in \mathcal{B} \) tal que \( y \in V \subseteq W \), lo que implica que \( \cl{V} \subseteq \cl{W} \subseteq Y \setminus \{z\} \implies z \notin \cl{V} \).

    Aplicando nuevamente la regularidad al abierto \( V \) que contiene a \( y \), existe un abierto un abierto \( O \) tal que \( y \in O \subseteq \cl{O} \subseteq V \).
    Tomando un elemento de la base \( U \in \mathcal{B} \) con \( y \in U \subseteq O \), obtenemos que \( \cl{U} \subseteq \cl{O} \subseteq V \).

    Esto demuestra que el par \( (U, V) \) pertenece a \( J \). Además, por construcción:
    \begin{itemize}
        \item \( y \in U \implies f_{U, V}(y) = 0 \).
        \item \( z \notin \cl{V} \implies z \in Y \setminus V \implies f_{U, V}(z) = 1 \).
    \end{itemize}
    Como \( f_{U, V}(y) \neq f_{U, V}(z) \), se tiene que \( F(y) \neq F(z) \), lo que prueba que \( F \) es inyectiva.

    2. \( F \) es un homeomorfismo sobre su imagen:
    Para ver que \( F : Y \to F(Y) \) es abierta, basta demostrar que para cualquier abierto básico \( W \in \mathcal{B} \) y cualquier \( y \in W \), existe un abierto \( H \) en la topología producto de \( {[0, 1]}^J \) tal que \( y \in F^{-1}(H) \subseteq W \).

    Dado \( y \in W \in \mathcal{B} \), por la regularidad de \( Y \) existen \( U, V \in \mathcal{B} \) tales que \( y \in U \subseteq \cl{U} \subseteq V \subseteq \cl{V} \subseteq W \). Así, \( (U, V) \in J \) y además \( f_{U, V}(y) = 0 \).

    Consideremos el abierto subbásico en el producto definido por:
    \[
        H = \{ x \in {[0, 1]}^J : x_{(U, V)} < 1 \}
    \]
    Su preimagen bajo \( F \) es \( F^{-1}(H) = \{ z \in Y : f_{U, V}(z) < 1 \} \).
    Si \( z \notin W \), entonces \( z \notin V \) (ya que \( V \subseteq W \)), lo que significa que \( z \in Y \setminus V \) y por lo tanto \( f_{U, V}(z) = 1 \). De este modo, si \( f_{U, V}(z) < 1 \), necesariamente \( z \in W \).

    Esto demuestra que \( y \in F^{-1}(H) \subseteq W \). Por lo tanto, la topología de \( Y \) coincide con la del subespacio \( F(Y) \), concluyendo que \( F \) es un embedding.
\end{proof}

\end{document}
